\section*{4. 结论}

结合定理~1.1 与命题~3.2 得到以下结果.

\subsubsection*{定理 4.1.}
固定特征 $p$、维数 $n$ 以及多重次数 $(d_1,\ldots,d_r)$. 记
\begin{equation}\tag{4.1.0}
X_V\longrightarrow V
\end{equation}
为特征 $p$ 时维数为 $n$、多重次数为 $(d_1,\ldots,d_r)$ 的泛光滑完全交. 若
\begin{equation}\tag{4.1.1}
\sum_{i=1}^{r}d_i\geq n+r+1,
\end{equation}
则存在稠密开子集 $U\subset V$, 使得
\begin{equation}\tag{4.1.2}
\text{对每个 }u\in U,\qquad
H^n(X_{\bar u},\mathbb Q_\ell)\ne N^1H^n(X_{\bar u},\mathbb Q_\ell).
\end{equation}

同样, 结合定理~1.3 与命题~3.2 得到以下结果.

\subsubsection*{定理 4.2.}
设 $X$ 是域 $k$ 上维数 $n\geq1$ 的光滑完全交. 固定整数
\begin{equation}\tag{4.2.1}
d\geq n+1.
\end{equation}
记
\begin{equation}\tag{4.2.2}
X\!\cdot\!H_u\longrightarrow U
\end{equation}
为 $X$ 的 $d$ 次超曲面截面中光滑且维数为 $n-1$ 者构成的泛族.
存在稠密开子集 $T\subset U$, 使得
\begin{equation}\tag{4.2.3}
\text{对每个 }t\in T,\qquad
H^{n-1}(X\!\cdot\!H_{\bar t},\mathbb Q_\ell)
\ne N^1H^{n-1}(X\!\cdot\!H_{\bar t},\mathbb Q_\ell).
\end{equation}

\subsection*{参考文献}

\begin{enumerate}[label={\textup{[\arabic*]}},leftmargin=2.6em]
\item J.-P. Serre, \emph{Faisceaux algébriques cohérents}, Annals of
  Mathematics \textbf{61} (1955), 197--278.
\item J.-P. Serre, \emph{Sur la topologie des variétés algébriques en
  caractéristique $p$}, in \emph{Symposium Internacional de Topología
  Algebraica}, Mexico, 1958.
\end{enumerate}

\medskip
\noindent\emph{底本注.}
照录法文层保留了印本 (4.1.1) 中的下限 $1$ 以及未标重音的短语 \emph{de meme}.
校订法文层和英文层将求和指标明确为 $i=1$, 并采用通常带重音的形式.
